\section{Limitations}

CVT assumes that the successor at an artificial cut is observed and belongs to the same continuing process. It must not bootstrap across a genuine reset, missing-data discontinuity, environment termination, or goal completion. The benchmark therefore applies only when boundary semantics can be identified well enough to distinguish continuation from true termination.

The primary study uses three optimization seeds and three segmentation realizations on PointMaze Medium. The independent validation adds a separate published $n$-step implementation on Puzzle-4x5, but it still uses only three optimization seeds and one fixed segmentation realization. Reported standard deviations are descriptive rather than well-powered inferential confidence intervals. These experiments support a cross-learner reliability failure, not a claim that every multi-step offline RL algorithm will exhibit the same magnitude of sensitivity. Broader validation should include additional algorithm families, manipulation domains with strong uncut performance, and naturally fragmented datasets with documented boundary causes.

CVT does not achieve perfect segmentation invariance in the primary study. At 100k updates it reaches 39.1\% versus 50.5\% for the uncut control, a residual gap of 11.4 percentage points. Approximation error, finite optimization, and the fact that horizon-indexed heads need not satisfy an exact continuing fixed point can all leave residual sensitivity even when boundary semantics are correct. Segmentation invariance should therefore be treated as a measurable reliability objective rather than an assumed binary property.
